\chapter{Contoh \textit{Caveat} Token Awal untuk Petugas Rekam Medis}
\label{lampiran:contoh-caveat-token-awal-rekam-medis}

Sebagai contoh, pasien dapat memberikan token awal kepada petugas rekam medis dengan cakupan akses yang cukup untuk mengoordinasikan pelayanan pada satu episode RME tertentu. Token tersebut dapat berisi \textit{caveat} sebagai berikut.

\begin{lstlisting}[caption={Contoh Caveat Token Awal untuk Petugas Rekam Medis}, label={lst:contoh-token-awal-rekam-medis}]
patient_address = 0xPASIEN
related_rme_id = RME-2026-000001
root_subject = 0xREKAM_MEDIS

read_dataset_in = [RAWAT_JALAN, RAWAT_INAP, LABORATORIUM, APOTEK]
write_dataset_in = [RAWAT_JALAN, RAWAT_INAP, LABORATORIUM, APOTEK]

read_function_in = [
  ANAMNESIS,
  PEMERIKSAAN_FISIK,
  DIAGNOSIS,
  TERAPI,
  PERMINTAAN_PEMERIKSAAN,
  HASIL_PEMERIKSAAN,
  DATA_RESEP_DAN_OBAT
  ...
]

write_function_in = [
  ANAMNESIS,
  PEMERIKSAAN_FISIK,
  DIAGNOSIS,
  TERAPI,
  PERMINTAAN_PEMERIKSAAN
  ...
]

expires_before = 2026-05-16T18:00:00
max_delegation_depth = 1
\end{lstlisting}

Nilai \texttt{0xPASIEN}, \texttt{RME-2026-000001}, dan \texttt{0xREKAM\_MEDIS} adalah contoh nilai yang hanya bersifat ilustratif. Pada kenyataannya, nilai pemegang token akan berbeda-beda tergantung pada alamat \textit{wallet} yang digunakan dan ID RME yang di-\textit{generate}. ID RME menggunakan format \texttt{RME-\{YYYY\}-\{sequence\_6\_digit\}} dengan nomor urut global per tahun pada \textit{Server} PRE. Simbol titik tiga (\texttt{...}) adalah bentuk ilustrasi yang digunakan untuk menunjukkan bahwa token dapat memiliki lebih banyak nilai lainnya selain yang dituliskan di sini.
